\subsection{Exact weights and approximants on the short mask}
On \(H_j\) the mass of the queens point \(p\) on each diagonal is at most \(C\alpha\), since its integer
length is less than \(\lceil\alpha j\rceil\). Turning off both diagonal factors changes the field in the
odds at each site by \(O(C\alpha+j^{-1})\). For the axial potentials we invert the margin Jacobian in the
centered gauge. Its Schur complement is \(I-p^{\mathsf T}p\) on centered columns. The gap bounds the
inverse there in the induced two-norm, and \(jp_i\le C\) bounds the sections of its normalized kernel in
Euclidean norm. The identity \((I-T)^{-1}=I+T+T(I-T)^{-1}T\) transfers the bound to the supremum norm for
\(T=p^{\mathsf T}p\). With \(A_{\rm ax}\) the axial incidence matrix and \(S_{\rm ax}=A_{\rm
ax}\operatorname{diag}(p)A_{\rm ax}^t\), the projector \(I-A_{\rm ax}^t S_{\rm ax}^+A_{\rm
ax}\operatorname{diag}(p)\) therefore has bounded supremum norm. For \(p^{\rm T}=p(1+z)\) we restrict
\(z\) to \(A_{\rm ax}\operatorname{diag}(p)z=0\). The derivative of \(\log p^{\rm T}+\log(1-p^{\rm T})\)
projected onto this space is \(I+O(\|z\|_\infty+j^{-1})\). Subtracting that gradient from \(z\) gives a
contraction on a ball of radius \(K(C\alpha+j^{-1})\) for a class constant \(K\). The resulting point is
the rooks maximum by rowwise concavity. If one entry satisfies \(t_1>1/2\), then for
\(\xi(t)=t(1-t)/(1-2t)\) we have \(\sum_{l>1}\xi(t_l)\le\xi(1-t_1)=-\xi(t_1)\) since \(\xi(t)/t\) is
increasing for \(t<1/2\). Cauchy--Schwarz with \(\sum_l z_l=0\) gives \(\sum_l z_l^2/\xi(t_l)\ge0\),
which is the Hessian of the negative action on the row. Strict local concavity identifies the branch. The
point \(p^{\rm T}\) satisfies \(p^{\rm T}_i/p_i=1+O(C\alpha+j^{-1})\). Keeping the margin \(p/2\le p^{\rm
T}\le3p/2\) gives \(j\max_i p^{\rm T}_i\le2C\), a gap \(\ge\min(1,\gamma)/2\) by Rayleigh, and
\(\epsilon_A:=\max_f\sum_{i\in f}p^{\rm T}_i\le2C\alpha\) in both diagonal families. The odds \(p^{\rm
T}_{rc}(1-p^{\rm T}_{rc})\) factorize into row and column factors on the support. Balancing with the
field \(\log(1-p^{\rm T}_{rc})=O(j^{-1})\) gives the exact doubly stochastic scaling \(D\) of the mask
\(A\), with \(D_i/p^{\rm T}_i=1+O(j^{-1})\). The absolute row and column sums of the perturbation are
\(O(j^{-1})\), and its operator norm has the same order.

\begin{proposicion}[Minors of the activation faces]\label{prop:menores-llegada}
Let \(U\subset A\) with \(v=|U|\le8\log n\). Remove the \(v\) sites to fix the exterior \(A_0=A\setminus
U\). For a face \(Y\subseteq U\), activate only its diagonal pairs. A subset \(I\subseteq U\) compatible
with the face and the axial constraints has weight \(W_I=\operatorname{per}(A_{0,-I})\), where \(-I\)
deletes the rows and columns of \(I\). On \(H_j\), if \(q=4v+12\), the order-\(q\) approximants
constructed below preserve the forbidden zeros and satisfy
\[
 |\log(\widehat W_I/W_I)|\le\eta_v=(Cq/j)^q+Ce^{-cj}.
\]
For \(5\le v\le8\log n\) and \(j\ge\sqrt n\), the summed error in the projections of the face logarithms
is \(O(j^{-8})\).
\end{proposicion}
\begin{proof}
Each placement determines a unique intersection \(I\) with \(U\). Deleting \(v\) entries perturbs the
margins by \(O(v/j)\), and deleting \(h=|I|\) rows and columns changes their normalized measures by
\(O(h/j)\) in total variation. The balancing inverse above gives each exterior its own scaling \(D'\),
including when the equation for an index of mass zero is extended. After identifying the spaces via the
square roots of the measures, the gap loss is \(O(\sqrt{h/j}+\sqrt v/j+(v+h)/j)\). The deleted sections
cost \(O(\sqrt{h/j})\) in Hilbert--Schmidt norm and the \(v\) entries cost \(O(\sqrt v/j)\). Thus
\(m=j-h\), \(mD'_{il}\le C'\), and \(B'=D'-\mathbf1\mathbf1^{\mathsf T}/m\) has operator norm \(\le
r'<1\). Hall's theorem applied to the support of \(D'\) gives \(W_I>0\).

Let \(R(z)=\sum_S z^{|S|}\prod_{e\in S}B'_e\), summed over partial matchings. For
\(T\sim\operatorname{Gamma}(m+1,1)\), the moments of \(y=m/T\) are \(m^k(m-k)!/m!\) up to \(k=m\).
Expanding the permanent gives \(m^m\operatorname{per}D'/m!=\mathbb E R(y)\). In a singular value
decomposition of \(B'\), take circular Gaussians \(Z_l\) of variance one and \(W=\sum_l\sigma_l|Z_l|^2\).
Wick's formula represents \(R(z)\) as the mean of \(\prod_i(1+\sqrt z\,u_i)\prod_k(1+\sqrt z\,v_k)\),
where \(u=U'\sqrt\Sigma Z,\ v=V'\sqrt\Sigma\overline Z\). Both fields sum to zero and their squared norms
are \(W\). The arithmetic--geometric mean inequality gives \(|R(z)|\le\mathbb E(1+|z|W/m)^m\). On
\(|z|r'<1\) with a positive margin, \(\log\mathbb E e^{|z|W}\le C\sum_l\sigma_l\le C\sqrt m\).

To define the truncation we extract the quadratic factor from the Wick expansion. This gives
\(d(y)=\det(I-y^2 B'^{\mathsf T}B')^{-1/2}\) and the pairings of the terms of \((1+t)e^{-t+t^2/2}\) of
degrees \(\ge3\), with covariance \(O(I+O)^{-1}\) and \(O=\left(\begin{smallmatrix}0&yB'\\yB'^{\mathsf
T}&0\end{smallmatrix}\right)\). With \(E\) edges and \(s\) active vertices the excess \(\ell=E-s\)
satisfies \(E\le3\ell\). The sum over assignments of distinct indices is expressed by inversion on
partitions of the \(s\) vertices. Blocks of repeated indices carry the Möbius factor \((-1)^{d-1}(d-1)!\)
for size \(d\). Each propagator is \(O(1/m)\) by diffusion and the resolvent with two exterior
sections. The number of pairings divided by \(s!\) is bounded by \(C^E E^{E-s}\). The coefficients at
each vertex and the lists of degrees cost \(C^E\). The absolute sum in the inversion, with one factor
\(m^{-1}\) per coincidence, costs \(\prod_{i<s}(1+i/m)\le e^{Cs}\) if \(s=O(m)\). Thus the contribution
\(V_\ell(y,1/m)\) to the remaining factor satisfies \(|V_\ell|\le(C\ell/m)^\ell\). We define
\(S_q=d(1+\sum_{\ell=1}^q V_\ell)\) as a function of \(y\) on a fixed disk about one.

Fix an auxiliary order \(Q_*=\lfloor\varepsilon j\rfloor\). In the variable \(w\) of \(R(yw)\), each edge
starts at degree one, so \(S_{Q_*}(yw)\) agrees through degree \(Q_*\). On \(|w|\le r_1>1\) with
\(r_1|y|r'<1\), both functions admit the majorant \(Ce^{C\sqrt m}\). At \(w=1\), Cauchy's estimate leaves
a remainder \(Ce^{C\sqrt m}r_1^{-Q_*}=O(e^{-cj})\). For \(y\) outside a fixed window about one, the Wick
bound cancels \(T^m\) in the gamma density. The substitution \(t'=t+W\) leaves the weight \(e^W\). If
\(W\le\kappa m\), the gamma tail remains at a distance proportional to \(m\) from its mean. On \(W>\kappa
m\) we use \(e^{-\xi\kappa m}\mathbb E e^{(1+\xi)W}\), choosing \((1+\xi)r'<1\). Both tails cost
\(O(e^{-cm+C\sqrt m})\).

For \(u=1/m\) the centered moments \(\mu_k=\mathbb E(y-1)^k\) satisfy
\[
 (1-ku)\mu_{k+1}=ku(2\mu_k+\mu_{k-1}),\qquad \mu_0=1,\quad\mu_1=0
\]
by integration by parts. In particular, \(\mu_{2h}\le(Ch/m)^h\) up to \(2h\le c m\) by induction. The sum
\(H_q=\sum_{k=0}^{2q+1}\mu_k\,\partial_y^k S_q(1,u)/k!\) integrates the Taylor expansion through degree
\(2q+1\). The next even moment and the fixed radius bound the remainder by \((C(q+1)/m)^{q+1}\), even
when the polynomial is extended outside the window. The sum of excesses from \(q+1\) to \(Q_*\) has the
same bound, since the ratio of majorants is \(\le Ce Q_*/m<1/2\). As \(H_q=d(1)+O(1/m)\), the logarithm
can be taken uniformly. For \(\widehat W_I\) we restore \(m!/m^m\) and all the exact scaling factors of
this minor. In particular, its normalizer uses \(m\). The positive comparison persists upon summing the
atoms of a face. Applying \(\Delta_U\) costs \(2^v\eta_v\). Since there are at most \(j^{2v}\) supports
and \(Cq\le j^{1/3}\) by the threshold \(n_0(b)\), the power term is \(\le2^v j^{-8-2v/3}\) per degree.
The total exponential term is \(\exp(O(\log^2 n)-cj)\). Both sum to \(O(j^{-8})\).
\end{proof}
